\begin{tikzpicture}[x=7.5mm, y=-8.5mm,
    zn/.style={komorka, minimum width=7.5mm, minimum height=7mm},
    znsz/.style={komorka szara, minimum width=7.5mm, minimum height=7mm},
    znok/.style={komorka ok, minimum width=7.5mm, minimum height=7mm},
    znzle/.style={komorka wyr, minimum width=7.5mm, minimum height=7mm}]
  \node[anchor=east, font=\small] at (-0.8,0) {\kod{S}};
  \foreach \z [count=\i from 0] in {b,a,n,a,n,y} {
    \node[zn] at (\i,0) {\z};
    \node[indeks] at (\i,-0.62) {\i};
  }
  \node[anchor=west, font=\small] at (6.6,0) {wzorzec \kod{W = "an"}, \ $n = 6$, $m = 2$};
  % kolejne położenia okna
  \foreach \i/\ok/\okno in {0/0/ba, 1/1/an, 2/0/na, 3/1/an, 4/0/ny} {
    \pgfmathsetmacro{\y}{\i+1.4}
    \node[anchor=east, font=\footnotesize] at (-0.8,\y) {$i = \i$};
    \foreach \z [count=\j from 0] in {b,a,n,a,n,y} {
      \pgfmathtruncatemacro{\w}{(\j>=\i) && (\j<=\i+1)}
      \ifnum\w=1
        \ifnum\ok=1 \node[znok] at (\j,\y) {\z}; \else \node[znzle] at (\j,\y) {\z}; \fi
      \else \node[znsz] at (\j,\y) {\z}; \fi
    }
    \ifnum\ok=1
      \node[anchor=west, font=\footnotesize] at (6.6,\y) {\kod{S[\i:\the\numexpr\i+2\relax]} = \kod{"\okno"} $=$ \kod{W} — \textcolor{zielen}{pasuje}};
    \else
      \node[anchor=west, font=\footnotesize] at (6.6,\y) {\kod{S[\i:\the\numexpr\i+2\relax]} = \kod{"\okno"} $\ne$ \kod{W}};
    \fi
  }
  \node[notka, anchor=west] at (6.6,6.6) {$n - m + 1 = 5$ położeń okna};
\end{tikzpicture}
